\subsection{A payment before the bill must be applied only once}
\label{sec:customer-advance}

The preceding recurrence deliberately excludes an advance before billing.
Suppose a customer pays 30 towards a service priced at 100. No service has yet
been performed and no payment was already due. The company now has cash 30 and
an obligation to provide the service, so it debits cash and credits a contract
liability. Crediting revenue would balance the entry but recognise performance
that has not occurred. Paragraphs 105--108 of the retained AASB 15 distinguish
this payment from an unconditional receivable \citep{aasb15dated}.

Now let the company earn 60 before issuing a bill. Recognition releases the
liability of 30 and creates a contract asset of 30, with a credit to revenue
of 60. The later invoice shows a gross amount of 80 and applies the customer's
advance of 30, leaving a new unconditional right to 50. Its entry debits the
receivable by 50, credits the contract asset by 30 and credits a contract
liability by 20. Creating a receivable of 80 would count the prepaid 30 again.
When the customer pays another 20, cash is 50, the receivable is 30 and the
contract liability is 20. Their net contribution to assets is the earned 60.

The method therefore keeps the invoice's gross amount separate from the right
still owed. Let $G$ be cumulative gross billings, $A$ cash advances, $K$ the
amount of those advances applied to bills, and $C$ later collections of
receivables. Each application names both its advance and its invoice; the
remaining unapplied credit is $A-K$. As before, $R$ is cumulative earned
revenue. The consideration already paid or made unconditionally due, counted
once, is $T=G+A-K$. For this one fixed-price contract,
\[
 \mathrm{AR}=G-K-C,\qquad
 \mathrm{CA}=\max(R-T,0),\qquad
 \mathrm{CL}=\max(T-R,0),\qquad
 \text{cash received}=A+C.
\]
The minus $K$ prevents the applied advance from appearing both in the cash
payment and in the bill. Since $\mathrm{CA}-\mathrm{CL}=R-T$, the control
\[
 (A+C)+(G-K-C)+\mathrm{CA}-\mathrm{CL}
 =G+A-K+R-T=R
\]
recovers the earned amount. This derivation uses a single contract, with
$0\leq K\leq\min(A,G)$, $0\leq C\leq G-K$ and $0\leq R,T\leq100$.
It provides no authority to offset different contracts or to ignore impairment,
financing, tax, cancellation or refund terms.

Each event must carry its contract, identity, date, currency and admitted
facts. In this restricted extension, advances precede all performance and
billing; later payments settle identified receivables. A revenue increment
first releases an existing contract liability and then creates a contract asset.
For a new gross bill $g$ that applies unused advance credit $k$, only $g-k$
creates a receivable. That amount releases the existing contract asset first;
any remainder increases the contract liability. The invoice consumes its
identified credit even if the advance has already been earned. A duplicate
credit application, an overcollection or missing unconditional-right evidence
cannot pass by adjusting a balancing account.

Billing before performance gives the same final balances in the example but
a different intermediate history. The invoice creates a receivable of 50 and
adds 50 to the existing contract liability of 30; earning 60 then releases 60
of that liability. After the later collection of 20, both paths have the
balances above. Completing the service adds revenue 40, releasing the remaining
liability of 20 and creating a contract asset of 20. A final bill for 20 and
collection of the remaining 50 leave cash 100 and no contract or receivable
balance. The actual sequence is retained even where endpoint amounts agree.

A wholly prepaid invoice exposes another implementation requirement. If its
gross amount and applied advance are both 30, it creates no new monetary right
and needs no accounting lines. It still consumes an identified credit. The
proposed authored method therefore declares a state-only journal proposal:
empty lines, zero account changes and explicit changes to the invoice and
advance allocations. Authorization checks that the net new right is zero;
the same atomic commit and replay rules protect the allocation from reuse.
An empty journal cannot carry a nonzero monetary effect. The native acquisition
schema requires nonempty lines and does not implement this case; an unavailable
authored schema returns an unsupported-operation result without changing state.

These are revenue and cash contributions before service costs and tax. Their
missing measurements remain inputs to the complete company statements. The
next example retains a different obligation: cash already received may have
to be returned even after the goods have been delivered.

\subsection{A return right survives the sale}
\label{sec:customer-returns}

Consider 100 goods sold for 10 each, with a former inventory cost of 6 each.
The buyer obtains control, pays cash and has a contractual right to a refund
for returned non-defective goods. Assume an eligible Australian for-profit
contract in 2024, within the retained AASB 15 compilation's application period.
The contract, control, price-estimation and significant-reversal conditions
are separately satisfied; immediate cash receipt alone would not establish
them. The numbers are synthetic measurements, not estimates of Apple's returns.

If the company initially expects 10 goods to be returned, it expects to retain
900 of the 1,000 received. AASB 15 B20--B25 requires revenue for the expected
entitlement, a refund liability and a separate asset for the right to recover
the goods \citep{aasb15dated}. With no expected recovery costs or deterioration
in this example, the refund liability is 100 and the recovery asset is 60.
The company debits cash 1,000, credits revenue 900 and credits the refund
liability 100. It removes inventory 600 by debiting cost of sales 540 and the
recovery asset 60. Pretax profit is therefore 360. Recognising the full 1,000
as revenue would omit a known refund obligation, although the entry balances.

For a homogeneous cohort of $N$ units, unit price $p$, former unit cost $c$ and
an admitted estimate $q$ of returns, the initial no-recovery-cost case gives
\begin{align*}
 \text{revenue}&=(N-q)p, & \text{refund liability}&=qp,\\
 \text{cost of sales}&=(N-q)c, & \text{recovery asset}&=qc.
\end{align*}
The estimate and its measurement date are retained with the cohort. Paragraphs
53--59 require an appropriate estimation method and the constraint on a
significant reversal; $q$ is not a probability supplied by the accounting
standard. Missing evidence for that measurement leaves the amount unresolved.
After actual returns begin, the state separately records refunded units,
remaining expected returns, remaining liability and recovery rights. A later
outcome cannot be inserted into the original information view.

Suppose the sale occurs in 2024 and new information at year end raises expected
returns to 12. The liability increases by 20, reducing revenue by 20. The
recovery asset increases by 12, reducing cost of sales by 12. Closing revenue
is 880, cost of sales 528 and pretax profit 352. The change in expectations
does not itself establish an error in the original estimate. If the opening
inventory of 600 was financed entirely by opening equity of 600, the resulting
pretax balance sheet has cash 1,000 and recovery asset 72, against refund
liability 120 and equity 952.

In 2025, assume the return period ends with eight accepted returns and no
remaining legal right or unresolved claim. Before final settlement, remeasure
the liability from 120 to 80 and the recovery asset from 72 to 48. The first
change credits revenue by 40; the second debits cost of sales by 24. Refunding
80 then debits the liability and credits cash. Receiving the eight goods
debits inventory by 48 and credits the recovery asset by 48. Settlement does
not recognise the sale or its cost a second time. The 2025 pretax profit
contribution is 16, closing cash is 920 and inventory is 48. With no other
movements, equity is 968: opening 600 plus cumulative profit 368.

The cash and accrual statements can now be checked independently. Operating
cash is 1,000 in 2024 and $-80$ in 2025. In the indirect method, the first
year adds the inventory reduction and liability increase to profit while
subtracting the recovery-asset increase. The second year reverses the relevant
balance movements:
\begin{align*}
 \mathrm{CFO}_{2024}&=352+600-72+120=1{,}000,\\
 \mathrm{CFO}_{2025}&=16-48+72-120=-80.
\end{align*}
There is no investing or financing cash flow, OCI or subsequent owner
transaction in this slice. Tax has not been computed: these pretax
contributions must enter the separate current/deferred-tax method before they
can become complete after-tax company statements.

Recovery costs change the asset without necessarily changing the refund. For
the 12-unit estimate, expected recovery costs of 1 per unit and deterioration
of 0.5 give a recovery asset of $12(6-1-0.5)=54$. Cost of sales becomes 546
and pretax profit 334, while the refund liability remains 120. The cost inputs
are measurements, not automatic zeros; estimating them does not pay cash early.
Paragraph B25 also requires separate presentation of the asset and liability.
Showing only their net 48 in the original no-cost case would conceal both
the gross right to recover goods and the gross obligation to refund cash.

Scope conditions matter as much as the arithmetic. Paragraph B26 treats an
exchange for the same type, quality, condition and price differently from a
return; B27 sends a defective-product replacement to warranty guidance.
The refund method cannot accept those events simply because the goods move
back to the seller. It also needs the contractual return window, accepted-unit
identities, actual refund and recovery measurements, and evidence that any
remaining rights have expired before clearing the balances. An unknown
estimate yields unresolved, absent control makes recognition ineligible, and
an event of a different kind remains outside this method. Inventory costing
and the condition of recovered goods continue to have their own requirements.
